\section{总结与延伸}

\subsection{最后的主张}

如果把整门 CS224N 的知识压成几条结论，大致是：

\begin{importantbox}{六条最后的 takeaway}
\begin{enumerate}
    \item \textbf{分布式表示是现代 NLP 的基础}：上下文和 dense vectors 改变了我们表示语言的方式。
    \item \textbf{Transformer 和预训练极其强大}：但强大不代表问题已经解决。
    \item \textbf{泛化、解释和多语言仍然开放}：这些不是边角，而是核心科学问题。
    \item \textbf{语言学依旧重要}：它帮助我们划分问题、分析失败和设计评测。
    \item \textbf{意义是渐变的、使用性的、组合性的混合体}：不能用单一理论一把抓。
    \item \textbf{能力增长伴随真实代价}：评测、对齐和治理必须和模型进步同步推进。
\end{enumerate}
\end{importantbox}

\subsection{一个更温和但更准确的结论}

人类语言不是单纯的符号，也不是单纯的统计模式。它既涉及世界中的对象，也涉及社会中的使用；既能被形式系统刻画，也会被大规模神经模型近似。CS224N 的最后一讲没有给出一个封闭答案，而是把问题推回到一个更诚实的位置：我们已经学会了做出很强的系统，但对语言、意义和智能本身，仍然需要更深的理论和更稳的实验。

\begin{figure}[H]
\centering
\includegraphics[width=0.95\textwidth]{slides-images/slide-064.jpg}
\caption{课程收尾页：回到 Stanford 的标志，结束在一个开放但不封闭的问题空间里\protect\footnotemark}
\end{figure}
\footnotetext{Slides 的结尾视觉上回到了 Stanford 的封印。}

\subsection{课程全局总结表}

\begin{table}[H]
\centering
\begin{tabular}{p{0.18\textwidth} p{0.22\textwidth} p{0.24\textwidth} p{0.23\textwidth}}
\toprule
主线问题 & 课程给出的答案 & 仍未解决的部分 & 对今天 NLP 的启示 \\
\midrule
语言怎么表示 & 分布式表示与上下文化 embedding 极其成功 & 表示是否对应可解释结构仍不清楚 & 表示学习仍是系统能力的底座 \\
语言怎么组合成意义 & Transformer 能学到大量可用组合规律 & 神经组合是否等价于人类式语义组合仍有争议 & 需要语言学和解释性研究补足 \\
模型是否真的泛化 & 在大规模任务上表现很强 & 数据效率、OOD 泛化、人类式学习仍不足 & benchmark 之外的评测要成为主线 \\
语言系统如何进入现实世界 & 预训练模型已能提升生产力 & 偏见、幻觉、能耗、版权与治理问题同步放大 & 系统设计必须把安全与社会成本前置 \\
\bottomrule
\end{tabular}
\caption{Lecture 18 的核心结论压缩表}
\end{table}

\subsection{给后续学习者的延伸路线}

\begin{enumerate}
    \item 从语言学和哲学角度继续追问 meaning、reference、use 和 compositionality
    \item 从 mechanistic interpretability 和 causal evaluation 角度研究模型内部到底学到了什么
    \item 从 multilingual NLP、speech、dialogue 和 grounded agents 的角度检验模型是否真的具备更强语言能力
\end{enumerate}

\subsection{如果重新学一遍 CS224N，建议这样重走}

\begin{table}[H]
\centering
\begin{tabular}{p{0.2\textwidth} p{0.23\textwidth} p{0.23\textwidth} p{0.22\textwidth}}
\toprule
阶段 & 应重点抓住的概念 & 最容易忽略的问题 & 重学时的关注点 \\
\midrule
表示学习阶段 & word vectors、distributional semantics、context & 把 embedding 当成纯技巧，不问它表示了什么 & 重新理解分布式表示与语义之间的关系 \\
序列建模阶段 & RNN/LSTM、attention、Transformer & 只记结构，不问 inductive bias 为何不同 & 对比不同架构的归纳偏置与数据效率 \\
预训练阶段 & language modeling、scaling、instruction tuning & 只看能力提升，忽略评测边界 & 把 pretraining 当作系统起点而不是终点 \\
系统与社会阶段 & reasoning、alignment、multilinguality、safety & 低估语言学和治理问题的重要性 & 把语言能力放进真实任务和社会环境里重新看 \\
\bottomrule
\end{tabular}
\caption{重学 CS224N 的建议路线}
\end{table}

这张表背后的意思很简单：技术章节最好和问题章节一起重读。否则很容易把 CS224N 误解成一门 ``神经网络模型年鉴''，而忽略它真正反复追问的那些核心问题。

\subsection{拓展阅读}

\begin{itemize}
    \item Manning, Christopher D. et al. CS224N 课程主页：\url{https://web.stanford.edu/class/cs224n/}
    \item Bender \& Koller, Climbing towards NLU: On Meaning, Form, and Understanding in the Age of Data：\url{https://aclanthology.org/2020.acl-main.463/}
    \item Ribeiro et al., "Why Should I Trust You?": Explaining the Predictions of Any Classifier：\url{https://arxiv.org/abs/1602.04938}
    \item Elhage et al., \emph{Transformer Circuits}：\url{https://transformer-circuits.pub/}
    \item Manning et al., The Stanford NLP Group: \url{https://nlp.stanford.edu/}
\end{itemize}

\end{document}
